\textbf{Note.} This chapter describes a \textbf{proposed operational design}. It is not implemented as a live system in this project.

\section{Deployment Objective}
\label{sec:8-1}

Provide recurring churn scoring that produces a prioritized existing-customer population for retention marketing.

\section{Weekly Batch Pipeline}
\label{sec:8-2}

Per-customer training weights ($1/n_i$) are \textbf{not} part of this pipeline; they exist only in training-data construction. Every active customer is scored every week.

\begin{center}
\textbf{Latest transactions} $\rightarrow$ \textbf{Valid transactions} $\rightarrow$ \textbf{Active customers} (Recency $\leq R_{\max}$) $\rightarrow$ \textbf{Remove suppressed customers} \\[0.3cm]
$\rightarrow$ \textbf{Snapshots} $\rightarrow$ \textbf{Features} $\rightarrow$ \textbf{Scoring} $\rightarrow$ \textbf{Ranking} $\rightarrow$ \textbf{Capacity \& business rules} $\rightarrow$ \textbf{Target population}
\end{center}

\section{Customer Risk Output}
\label{sec:8-3}

Minimum output: \texttt{customer\_id}, \texttt{snapshot\_date}, \texttt{churn\_probability}, \texttt{churn\_probability\_calibrated} (Section~\ref{sec:6-quality}), \texttt{expected\_loss} (calibrated probability $\times$ $V$ $\times$ expected orders in $H$, Section~\ref{sec:7-sens}), \texttt{risk\_band} (fixed cut-offs on the calibrated probability, Section~\ref{sec:6-quality}), \texttt{rank} (by expected loss within the week, 1 = first to contact), \texttt{segment}, \texttt{model\_version}.

\section{Target Population}
\label{sec:8-4}

Select target customers based on churn risk, customer value, marketing capacity and business constraints. The recommended ranking key is \texttt{expected\_loss}, not \texttt{churn\_probability}: under a realistic, suppressed campaign with a small weekly budget, ranking by expected loss preserves the most net value of every candidate evaluated, while ranking by churn probability alone does worse than the recency rule (Section~\ref{sec:7-10}). The weekly budget should stay small relative to the active base -- with a 13-week suppression, a 10\% weekly share contacts most customers within a quarter, after which the choice of ranking barely matters.

\section{Suppression}
\label{sec:8-5}

Customers \textbf{contacted} in a campaign are excluded from the Scoring Population for a configurable suppression period $S$, so they are not targeted repeatedly. The evaluation uses $S = 13$ weeks ($\approx H$, one contact per churn window; Section~\ref{sec:7-sens}); shorter periods (e.g.\ 4--8 weeks) allow follow-up contacts within the same churn window. Suppression:

\begin{itemize}
    \item applies only to customers who were actually targeted, not to everyone who was scored;
    \item is an operational business rule, independent of how training rows are sampled or weighted;
    \item is logged, so that suppressed customers can be analyzed separately once campaign data exist.
\end{itemize}

\section{Monitoring}
\label{sec:8-6}

Data quality, feature drift and business metrics where available, plus the weekly model checks of Section~\ref{sec:6-quality} against the validation reference. They split by when they become available:

\begin{itemize}
    \item \textbf{Immediately, without labels:} HHI and CSI of the risk-band distribution. CSI is the early warning -- on the test period it already crossed 0.10 within six weeks as the churn rate fell.
    \item \textbf{After $H = 90$ days, once labels mature:} AUC, Brier score, and the mean calibrated probability against the observed churn rate.
\end{itemize}

Calibration is tied to the churn rate of the data it was fitted on, so it drifts whenever that rate moves -- as it does seasonally here. Two responses, from light to heavy:

\begin{itemize}
    \item \textbf{Recalibrate} (quarterly, or when the mean calibrated probability and the observed churn rate differ by more than about 5 percentage points): refit only the log-odds offset $\delta$ on the latest weeks with matured labels. The model, its ranking and the band order are untouched; the band cut-offs are recomputed from the same weeks. Because labels arrive $H$ days late, the recalibrated rate always reflects the previous quarter -- the churn rate of the same season a year earlier is a reasonable alternative target.
    \item \textbf{Retrain} when AUC falls clearly below its validation level or CSI exceeds 0.25, or whenever the churn definition itself changes (e.g.\ a different $H$), which invalidates both the model and its calibration.
\end{itemize}
